\section{The Fourier window of the actual curvature covariance}
\label{sec:actual-covariance-fourier-window}

The Gaussian comparison requires estimates for the covariance selected
by the construction, uniformly in its ancestry and physical period.
We derive these estimates before considering the cumulative transport.
Fix the positive kernel penalty and a smoothing order $J\ge4$.
All constants in this section may depend on these two choices and on
the finite derivative order, but not on block lengths or ancestry.

\subsection{A uniformly elliptic symbol}

Let $H(p)$ be the two-form pulse symbol of the chosen ancestry, in
radian frequency $p\in\mathbb R^4$. The component and positive-row
constructions give constants $\alpha,\eta>0$ and $\Omega<\infty$
such that its Laurent coefficients have strip cost $\Omega$ at
width $\alpha$, its action on every finite period has quadratic
lower bound $\eta$, and $H(0)=I$. These are properties of the
original pulse, before any replacement of the curvature completion.
Exponential summability implies, for every integer $m\ge0$,
\[
 \sup_p\|D^mH(p)\|\le C_m.
\]
Indeed each derivative multiplies the coefficient at $z$ by at most
$|z|^m$, and $|z|^m e^{-\alpha|z|_1}$ is bounded. The finite-period
lower bound also holds at every real frequency. First extend the
quadratic bound of the real finite-period operator to complex vectors
by splitting them into real and imaginary parts. Then test
the period-$n$ matrix on the vector $e^{ip\cdot x}v$ for a grid
frequency $p$. Its quadratic form and squared norm both acquire the
same factor $n^4$, which cancels. Approximate an arbitrary real
frequency by grid frequencies and use continuity. Thus
\[
       \operatorname{Re}\langle H(p)v,v\rangle\ge\eta|v|^2.
\]

Put $q_a(p)=e^{ip_a}-1$, and let $X(p)$ and $U(p)$ be exterior
multiplication by $q(p)$ from one-forms to two-forms and from
scalars to one-forms, respectively. Define
\[
 F(p)=X(p)^*H(p)X(p)+U(p)U(p)^*,\qquad
 P(p)=X(p)F(p)^{-1}X(p)^*\quad(p\ne0),
\]
and set $P(0)=0$. On the cube $|p_a|<3\pi/2$, the elementary
identity for exterior multiplication and contraction gives
\[
 \|Xv\|^2+\|U^*v\|^2=|q|^2\|v\|^2,
 \qquad
 \operatorname{Re}\langle Fv,v\rangle
       \ge\min(\eta,1)|q|^2\|v\|^2.
\]
Here $|q|\ge c|p|$: the continuous function
$|e^{it}-1|/|t|$ has a positive lower bound on
$[-3\pi/2,3\pi/2]$, with its value at zero defined to be one.
Consequently $F(p)$ is invertible for $p\ne0$ in this cube.

We record the derivative estimates with their exact gain. A function
has order $s$ if, through the finite number of derivatives under
consideration, its $m$th derivative is bounded by
$C_m|p|^{s-m}$. Products add orders by Leibniz' rule.
Taylor's formula gives $q-ip$ order two and $H-I$ order one;
$q$ itself has order one. The exterior identity at $q=ip$ gives
\[
 F(p)=|p|^2I+E(p),\qquad E\text{ has order }3.
\]
Thus $T=|p|^{-2}F$ has order zero, $T-I$ order one, and
$\|T^{-1}\|$ is uniformly bounded by coercivity. Differentiating
$TT^{-1}=I$ proves inductively that $D^mT^{-1}=O(|p|^{-m})$:
the term containing $D^mT^{-1}$ is isolated and every remaining
term is a product of derivatives whose orders total $m$.
Furthermore $T^{-1}-I=T^{-1}(I-T)$ has order one. It follows that
\begin{equation}\label{eq:actual-symbol-orders}
 P\text{ has order }0,\qquad P-P_0\text{ has order }1,
 \quad P_0(p)=X_0(p)|p|^{-2}X_0(p)^*,
\end{equation}
where $X_0$ is exterior multiplication by $ip$.
These estimates are uniform over all the selected ancestries.

We verify that $P$ is the actual curvature multiplier. Let $D$ be
the original one-form covariance and $R=X^*HX$. The transport
identities already proved give $RDX^*=X^*$ and $XU=0$.
Write $B=UU^*$. Then $XB=0$, $BX^*=0$, and $RB=BR=0$.
Hence $F=R+B$ commutes with $B$, as does its inverse, and
$XF^{-1}B=0$. Multiplying $RDX^*=X^*$ on the left by
$XF^{-1}$ yields
\[
 XF^{-1}X^*=XF^{-1}RDX^*=XDX^*.
\]
At zero both curvature multipliers vanish. Thus no change of the
finite-period covariance has been made.

\subsection{A smooth chart without changing the finite Fourier sum}

Choose a smooth nondecreasing function $g$ that is zero on
$(-\infty,-5\pi/4]$ and one on $[-3\pi/4,\infty)$, and put
\[
 h(p)=\prod_{a=1}^4\bigl(g(p_a)-g(p_a-2\pi)\bigr).
\]
It is supported in $|p_a|\le5\pi/4$, equals one near zero,
and satisfies
\[
                 \sum_{z\in\mathbb Z^4}h(p+2\pi z)=1.
\]
The last identity is the product of four telescoping sums; only
finitely many terms are nonzero. Define the unit-frequency chart
$K(\xi)=h(2\pi\xi)P(2\pi\xi)$, entry by entry and with the
usual Kronecker delta on colour indices. Its support lies in
$|\xi|\le2$. For an integer displacement $x$, periodicity and
the partition identity show that the exact period-$n$ covariance is
\begin{equation}\label{eq:actual-chart-sum}
 \kappa_n(x)=n^{-4}\sum_{z\in\mathbb Z^4}
              e^{2\pi i x\cdot z/n}K(z/n).
\end{equation}
The sum is finite. Grouping $z$ by its residue modulo $n$ proves
the formula, including even $n$ and the zero frequency. At zero
we retain the actual value $P(0)=0$, rather than assigning a
continuous value to a direction-dependent limit.

For two lists of coordinate directions, of lengths $a,b$, let
$u=a+b$. Differencing the first endpoint multiplies $K(\xi)$
by $e^{2\pi i\xi_c}-1$; differencing the second endpoint uses
$e^{-2\pi i\xi_c}-1$. Denote their product by $m(\xi)$,
and the corresponding product of $2\pi i\xi_c$ and
$-2\pi i\xi_c$ by $m_0(\xi)$. Repeatedly using
$e^{it}-1-it=O(t^2)$, with its differentiated Taylor bounds,
shows that $m-m_0$ has order $u+1$ and $m,m_0$ have order $u$.
By \eqref{eq:actual-symbol-orders}, the remainder
\[
 G(\xi)=m(\xi)K(\xi)
       -m_0(\xi)e^{-(2\pi)^2|\xi|^2}P_0(2\pi\xi)
\]
satisfies, for $0\le v\le u+6$,
\begin{align}\label{eq:actual-mist}
 \|D^vG(\xi)\|&\le C|\xi|^{u+1-v}
                              &&(0<|\xi|\le1),\\
 \|D^vG(\xi)\|&\le C|\xi|^{-5}
                              &&(|\xi|\ge1).\notag
\end{align}
For the first estimate, subtract and add $m_0K$ and use the
product rule. The extra Gaussian differs from one to order two.
Outside the unit ball the compactly supported term has bounded
derivatives, and every derivative of the Gaussian term decays
faster than any power. Finite maxima give the same $C$ for all
face pairs and direction lists of bounded length.

\subsection{Fourier decay and quadrature at the pole}

Write $\mathcal F_+f(x)=\int e^{2\pi i x\cdot\xi}f(\xi)\,d\xi$.
We first show directly that \eqref{eq:actual-mist} implies
\begin{equation}\label{eq:actual-mist-decay}
             |\mathcal F_+G(x)|\le C'|x|^{-u-5}
                                      \quad(|x|\ge1).
\end{equation}
Set $r=|x|^{-1}$ and choose a smooth radial cutoff equal to zero
on the ball of radius $r/2$ and one outside the ball of radius
$r$. The omitted inner term has integral at most
$C\int_0^r t^{u+1}t^3\,dt\le C'r^{u+5}$.
For the outer term integrate by parts $u+6$ times in the unit
direction $x/|x|$. Each term where a derivative hits the cutoff
is supported on $r/2\le|\xi|\le r$ and has integral
$O(Cr^{-1})$. For the remaining terms, derivatives of order
$u+6$ have integral bounded by
\[
 C'\left(\int_r^1 t^{-5}t^3\,dt+
                    \int_1^\infty t^{-5}t^3\,dt\right)
                  \le C''/r.
\]
Boundary terms vanish by first inserting a distant smooth cutoff
and then sending its radius to infinity; the displayed bounds
control its derivatives as well. Multiplication by
$|x|^{-u-6}$ proves \eqref{eq:actual-mist-decay}.

Next let $f=mK$. It is compactly supported, bounded by $C$, and
satisfies $\|Df(\xi)\|\le C/|\xi|$ off zero. Partition frequency
space into cubes of side $1/n$, and use their grid corner as the
sample point. Its distance from any point of that cube is at most
$2/n$. If this distance is less than $|\xi|/2$, the mean value
bound gives
\[
 |f(\xi)-f(z/n)|\le C'/(n|\xi|).
\]
Otherwise the same estimate follows from the bound $2C$ on the
difference, after increasing $C'$. This also covers a sampled
corner at zero. The phase difference is at most
$2\pi|x|\,|\xi-z/n|$. Integration over a fixed enlarged support
ball is legitimate because $|\xi|^{-1}$ is locally integrable in
four dimensions. We obtain
\begin{equation}\label{eq:actual-quadrature}
 \left|\mathcal F_+f(x)-n^{-4}\sum_z
                e^{2\pi i x\cdot z/n}f(z/n)\right|
                  \le C'\frac{1+|x|}{n}.
\end{equation}
This estimate treats the singular grid point without removing it.

Finally the Gaussian-regularized continuum term has the required
Maxwell limit. In the present Fourier convention its scalar
potential multiplier, for $t>0$ (here $t=(2\pi)^2$), is
$(2\pi)^{-2}e^{-t|\xi|^2}/|\xi|^2$. The identity
\[
 \frac{e^{-t|\xi|^2}}{|\xi|^2}
       =\int_t^\infty e^{-s|\xi|^2}\,ds
\]
and the product Gaussian integral give, at $x\ne0$,
\[
 \mathcal F_+\left((2\pi)^{-2}
          \frac{e^{-t|\xi|^2}}{|\xi|^2}\right)(x)
 =\frac{1}{4\pi^2|x|^2}
                 \left(1-e^{-\pi^2|x|^2/t}\right).
\]
Fubini is justified by
$\int_t^\infty(\pi/s)^2\,ds<\infty$.
All polynomial frequency moments are integrable, so differentiation
under the integral supplies both curls and both endpoint derivative
lists, with the opposite sign at the second endpoint. The difference
from the derivatives of $1/(4\pi^2|x|^2)$ is a derivative of a
Gaussian times that potential. It is bounded by any prescribed
inverse power for $|x|\ge1$.

Combining this fact with \eqref{eq:actual-chart-sum},
\eqref{eq:actual-mist-decay} and \eqref{eq:actual-quadrature} proves
\begin{equation}\label{eq:actual-native-model}
 \bigl|\Delta_x^a\Delta_y^b\kappa_n(x,y)
                     -\mathcal K_{a,b}(x-y)\bigr|
 \le C_{a,b}\left(|x-y|^{-a-b-5}
                          +\frac{1+|x-y|}{n}\right),
 \qquad |x-y|\ge1.
\end{equation}
Here each superscript denotes its specified direction list, and
$\mathcal K_{a,b}$ is the corresponding continuum Maxwell curvature
covariance. The constants are uniform in ancestry and period.
The covariance has closed rows and columns because its construction
is $dDd^*$ and $d^2=0$. Its entries are real by conjugating the
finite Fourier sum and replacing the frequency by its negative.
No positivity of a later normal-ordered covariance is inferred.

For a fixed comparison order $d\ge1$, only $a,b\le d-1$ are needed.
The zero-order comparison can use $d=1$ and discard the extra data.
The construction reserves their finitely many chart constants at
smoothing order $J$ before choosing the block length. At the chosen
comparison level the reserved chart cost is paid by $H^{1/4}$,
where $H$ is its monitor. The finite constants from the multiplier
and Fourier estimates are then absorbed successively into
$H^{1/2}$ and $H$. This yields \eqref{eq:actual-native-model}
with $C_{a,b}$ replaced by $H$, on that same chosen history.
The quantitative clearance at the chosen level remains part of
the common scale selection; merely taking an unrelated level large
would not establish it.
